\documentclass[12pt,a4paper]{article}
\usepackage[utf8]{inputenc}
\usepackage[T1]{fontenc}
\usepackage{geometry}
\usepackage{listings}
\usepackage{xcolor}
\usepackage{hyperref}

\geometry{margin=1in}

\definecolor{codegreen}{rgb}{0,0.6,0}
\definecolor{codegray}{rgb}{0.5,0.5,0.5}
\definecolor{codepurple}{rgb}{0.58,0,0.82}
\definecolor{backcolour}{rgb}{0.95,0.95,0.92}

\lstdefinestyle{mystyle}{
    backgroundcolor=\color{backcolour},   
    commentstyle=\color{codegreen},
    keywordstyle=\color{magenta},
    numberstyle=\tiny\color{codegray},
    stringstyle=\color{codepurple},
    basicstyle=\ttfamily\footnotesize,
    breakatwhitespace=false,         
    breaklines=true,                 
    captionpos=b,                    
    keepspaces=true,                 
    numbers=left,                    
    numbersep=5pt,                  
    showspaces=false,                
    showstringspaces=false,
    showtabs=false,                  
    tabsize=2
}

\lstset{style=mystyle}

\title{JEE Project: JSP and Servlet Implementation Report}
\author{Gemini CLI Assistant}
\date{May 21, 2026}

\begin{document}

\maketitle

\section{Introduction}
This report provides a theoretical overview of how the Student Management System would be implemented using JavaServer Pages (JSP) and Servlets. This approach leverages server-side rendering, as opposed to the client-side rendering used in the JavaFX/REST architecture.

\section{JSP/Servlet Architecture}
In a JSP-based architecture, the server handles both the business logic and the generation of the user interface (HTML).
\begin{itemize}
    \item \textbf{Controller (Servlets)}: Intercepts HTTP requests, interacts with repositories, and forwards data to JSPs.
    \item \textbf{View (JSP)}: Dynamically generates HTML using data provided by the servlets.
    \item \textbf{Model}: Reuses the existing JPA entities (\texttt{User}, \texttt{Student}).
\end{itemize}

\section{Core Components}

\subsection{Authentication Servlet}
\textbf{Description}: Handles login submissions via POST.
\begin{itemize}
    \item \textbf{Servlet Name}: \texttt{LoginServlet}
    \item \textbf{URL Pattern}: \texttt{/login}
    \item \textbf{Process}:
    \begin{enumerate}
        \item Retrieves \texttt{username} and \texttt{password} parameters from the request.
        \item Injects \texttt{UserRepository} and \texttt{PasswordService}.
        \item Validates credentials.
        \item On success: Stores user info in \texttt{HttpSession} and redirects to \texttt{dashboard.jsp}.
        \item On failure: Sets an error attribute and forwards back to \texttt{login.jsp}.
    \end{enumerate}
\end{itemize}

\subsection{Student Management Servlet}
\textbf{Description}: Manages CRUD operations for students.
\begin{itemize}
    \item \textbf{Servlet Name}: \texttt{StudentServlet}
    \item \textbf{URL Pattern}: \texttt{/students}
    \item \textbf{Actions}:
    \begin{itemize}
        \item \texttt{GET}: Fetches all students and forwards to \texttt{student-list.jsp}.
        \item \texttt{POST} (Create): Parses form data and calls \texttt{studentRepository.save()}.
        \item \texttt{POST} (Delete): Uses a hidden \texttt{\_method} parameter or specific action parameter to call \texttt{studentRepository.delete()}.
    \end{itemize}
\end{itemize}

\section{View Implementation (JSP)}

\subsection{student-list.jsp}
This page uses JSTL (\texttt{c:forEach}) to iterate over the student collection:
\begin{lstlisting}[language=HTML]
<table>
    <c:forEach items="" var="s">
        <tr>
            <td></td>
            <td></td>
            <td></td>
        </tr>
    </c:forEach>
</table>
\end{lstlisting}

\section{Conclusion}
Transitioning the current REST backend to a JSP-based system would involve replacing the JAX-RS resources with Servlets and creating JSP templates for the UI. The core business logic and database layers (JPA and Repositories) remain completely reusable, demonstrating the flexibility of the JEE architecture.

\end{document}
